\documentclass{article}
\usepackage{etoolbox}
\usepackage{graphicx}
\usepackage{amsmath}

\ifdef{\myoutput}{}{\def\myoutput{concept}}
\ifdef{\myversion}{}{\def\myversion{1}}

\usepackage[output=\myoutput
,numberofversions=4
,version=\myversion
,seed=4
,randomizequestions=true
,randomizeanswers=true
,writeRfile=true
]{mcexam}

\usepackage[margin=0.95in]{geometry}
\usepackage{multicol}

\usepackage{fancyhdr,lastpage}
\pagestyle{fancy}
\fancyhf{}
\renewcommand{\headrulewidth}{1pt} 
\renewcommand{\footrulewidth}{1pt}
\lfoot{\mctheversion}
\rfoot{Page \thepage\ of \pageref{LastPage}}
\lhead{Quiz 9.4 \mctheversion \\ MAT 2000: Introduction to Statistics}
\rhead{Name Print): \hspace{2.5 in}}

\usepackage{calc}
\renewenvironment{setmcquestion}
{\begin{minipage}[t]{\linewidth-\labelwidth}}
	{\end{minipage}\par}
	
\renewenvironment{setmcquestioninstruction}
{\begin{minipage}{\textwidth}}
	{\end{minipage}}
	
\renewenvironment{setmcanswers}{}{}
\setlist[setmcquestions]{label=\mcquestionlabelfmt{*}.
	,ref=\mcquestionlabelfmt{*}
	,itemsep=0.5\baselineskip
	,topsep=1\baselineskip
}

\renewcommand\mcversionlabelfmt[1]{\arabic{#1}}

%copy-pasting from MS Word, use regex to find [\*]*[a-d]\. and replace with \answer 

\raggedbottom

\begin{document}
	Circle the correct answer for each question.  If you're unsure, you can write numbers 1, 2, 3, etc. next to the choices to get partial credit.  Questions refer to the nonindented paragraph above unless otherwise stated.\\[.5cm]
	
	\begin{mcquestions}
		
		\begin{mcquestioninstruction}
			According to the 2019 U.S. Census Bureau estimates, 63.4\% of the U.S. self-identifies as Non-Hispanic white, 15.3\% as Hispanic and Latino, 13.4\% as Black or African American, 5.9\% as Asian, and 2.0\% as Other. A random sample of 1000 students graduating from California colleges and universities resulted in the following data on self-identified race:
			\begin{center}
				\begin{tabular}{| p{1in} | p{1in} | p{1in} | p{1in} | p{1in} | p{1in} |}
					\hline
					\textbf{Self-identified Race} & Non-Hispanic White & Black or African American & Hispanic or Latino & Asian & Other \\
					\hline
					\textbf{Number in Sample} & 679 & 51 & 77 & 190 & 3 \\
					\hline
				\end{tabular}
			\end{center}
			Use this information to answer questions \ref{census-1} and \ref{census-2}
		\end{mcquestioninstruction}
		
		\question In the null hypothesis, what does the symbol \(\mu_{\text{Race}}\) represent? \label{census-1}
		\begin{mcanswerslist}
			\answer  The proportion of the U.S. that self-identify as a certain Race
			\answer  The proportion of Californians that self-identify as a certain Race
			\answer[correct]  The proportion of all students graduating from California colleges and universities that self-identify as a certain Race
			\answer  The proportion of the sample of 1000 students that self-identified as a certain Race
		\end{mcanswerslist}
		
		\question[follow] Based upon a p-value of less than 0.001, what is an appropriate conclusion for the goodness-of-fit test? \label{census-2}
		\begin{mcanswerslist}
			\answer  Strong evidence does not exist that the distribution of self-identified race among California college and university students differs from the distribution of self-identified race in the U.S. overall.
			\answer[correct]  Strong evidence exists that the distribution of self-identified race among California college and university students differs from the distribution of self-identified race in the U.S. overall.
			\answer  Strong evidence exists that the distribution of self-identified race among California college and university students is the same as the distribution of self-identified race in the U.S. overall.
			\answer  Strong evidence exists that the probabilities of each self-identified race among California college and university students are not all equal.
		\end{mcanswerslist}
		
		\begin{mcquestioninstruction}
			Are people equally likely to be born on any day of the seven days of the week? Or are some days more likely to be a person's birthday than other days?  To investigate this question, days of birth were recorded for the 147 "noted writers of the present" listed in The World Almanac and Book of Facts 2000. The counts for the seven days of the week are given in the following table.
			\begin{center}
				\begin{tabular}{ccccccc}
					Mon&Tues&Wed&Thurs&Fri&Sat&Sun\\ \hline
					17&26&22&23&18&15&25					
				\end{tabular}
			\end{center}
			Use this information to answer questions \ref{birthday-1} and \ref{birthday-2}
		\end{mcquestioninstruction}
		\question Are validity conditions met to conduct a theory-based chi-square goodness-of-fit test? \label{birthday-1}
		\begin{mcanswerslist}
			\answer[correct]  Yes, since all observed counts are at least 10.
			\answer  Yes, since we have at least 10 successes and 10 failures.
			\answer  No, since the p-values in questions 35 and 37 were different.
			\answer  No, since this was not a random sample.
		\end{mcanswerslist}
		
		\question[follow] Based upon a p-value of 0.53, what is the appropriate conclusion for the test? \label{birthday-2}
		\begin{mcanswerslist}
			\answer  Strong evidence exists that people are equally likely to be born on any day of the seven days of the week.
			\answer  Strong evidence does not exist that people are equally likely to be born on any day of the seven days of the week.
			\answer  Strong evidence exists that people are not equally likely to be born on any day of the seven days of the week.
			\answer[correct]  Strong evidence does not exist that people are not equally likely to be born on any day of the seven days of the week.
		\end{mcanswerslist}
		
		
	\end{mcquestions}
\end{document}